\documentclass[11pt]{article}
\usepackage[margin=2.4cm]{geometry}
\usepackage{fontspec}
\setmainfont{DejaVu Sans}
\usepackage{juliamono}
\usepackage{fvextra}
\usepackage{xcolor}
\usepackage{xurl}
\usepackage{hyperref}
\usepackage{amssymb}
\newcommand{\checkbox}{$\square$}
\newcommand{\answerchoices}{\checkbox\ YES \quad \checkbox\ NO \quad \checkbox\ CANNOT DETERMINE}
\newcommand{\decisionexplanation}{%
  \par\smallskip\textbf{If NO, give the reason for the mismatch (required). If
  CANNOT DETERMINE, state what context or evidence is missing:}\par
  \rule{\linewidth}{0.4pt}\par\vspace{7pt}\rule{\linewidth}{0.4pt}\par}
\setlength{\parindent}{0pt}\setlength{\parskip}{5pt}
\setlength{\emergencystretch}{2em}
\DefineVerbatimEnvironment{Vsrc}{Verbatim}{fontsize=\small,frame=single,
  breaklines=true,breakanywhere=true}
\DefineVerbatimEnvironment{Vlean}{Verbatim}{fontsize=\scriptsize,frame=single,
  numbers=left,numbersep=5pt,breaklines=true,breakanywhere=true}
\hypersetup{colorlinks=true,urlcolor=blue}

\begin{document}
\begin{center}{\LARGE Proof Verification Document}\\[2pt]{\large SE\_math\_4219732\_\_q}\\[2pt]{\small Stack Exchange, question arm (confirmed)}\end{center}
\begin{center}\fcolorbox{blue!60!black}{blue!4}{\parbox{0.92\linewidth}{\textbf{Make three independent decisions.} Follow the evidence chain: original source $\longrightarrow$ challenged claim $\longrightarrow$ proposed counterexample $\longrightarrow$ Lean theorem. A record is valid only when all three links are correct. Complete the human sign-off before reading the automated review appendix; those earlier results are background, not an answer key.}}\end{center}
\textbf{Use only this document, except for the public source link.} Everything needed for Questions 2 and 3 is included here. For Question 1, open the public source link when one is shown. If there is no public link, the complete stored problem and solution are included here. If the information is still not enough, choose CANNOT DETERMINE.
For each question, choose \textbf{YES}, \textbf{NO}, or \textbf{CANNOT DETERMINE}. Choose CANNOT DETERMINE if the document does not provide enough information for you to decide; do not mark YES if you are uncertain. If you find a specific mismatch, choose NO and explain it in the space below that question. Do not rewrite the claim, add assumptions, or substitute a different counterexample.
\textbf{How copied text is shown.} The problem, quoted passage, and proposed counterexample come from the stored record. The Lean theorem statement and Lean files are copied exactly from the certificate. Some long lines continue on the next printed line; this is only page formatting and does not change the copied text. Line numbers are added only to help you find passages. If the stored record contains a character that cannot be printed, the document shows \texttt{[U+XXXX CONTROL]} in its place.
\begin{description}
\item[Challenged claim.] The exact statement from the source that this record says is false.
\item[Proposed counterexample.] A concrete case that should satisfy the claim's assumptions but make its conclusion false.
\item[Lean certificate.] A Lean file proving a formal theorem. Successful compilation proves the theorem as written, not that it is the right translation of the source.
\end{description}
\section*{Question 1 --- Does this proof verification document accurately represent the source?}
\textbf{What to decide.} Compare the challenged passage below with the original source---either the public link or the complete source text in this document. Answer YES only if they make the same claim and include the same assumptions or conditions. If this document changes the meaning, leaves out a condition, or points to a different passage, answer NO.
\subsection*{Original source}
\textbf{Source record:}
\begin{Vsrc}
qedbench_certificate_availability_triage/data/records_se.jsonl
\end{Vsrc}
\textbf{Public source:}\par
\url{https://math.stackexchange.com/questions/4219732/rewriting-an-equation-in-the-terms-of-algorithmic-complexity}\par
\subsection*{Problem or question from the source (verbatim)}
\begin{Vsrc}
Rewriting an equation in the terms of Algorithmic complexity.
\end{Vsrc}
\subsection*{Challenged passage recorded from the source (stored text)}
\begin{Vsrc}
$$\binom{n}{2b+1}\left( \frac{2}{cn\cdot 2^f}\right)^{2b}= \frac{n^{2b+1}}{(2b+1)^{2b+1}} \frac{(2b+1)^{2b+1}}{2^f\cdot cn} = \frac{n}{1}\cdot\frac{1}{4^{f}\cdot n}=\frac{1}{4f}$$ Note that $n^{2b+1}=n^{2b}\cdot n = 2n = n$
\end{Vsrc}
\subsection*{Surrounding source lines (stored text with added line labels)}
The labels refer to lines in the solution text stored with the source record; the source-line content is not transliterated.
\begin{Vsrc}
0001 | In [this](https://www.cs.cmu.edu/~dga/papers/cuckoo-conext2014.pdf) paper (page 7 Eq. 3) the author derives the following: $$\binom{n}{2b+1}\left( \frac{2}{cn\cdot 2^f}\right)^{2b}=\Omega\left( \frac{n}{4^{bf}}\right)$$
0002 | I have tried to do the same below, but first some preliminaries. This is an analysis for algorithms and therefore constants can be erased. In this example, $b$ is a constant while $f$ is some $f$-bit fingerprint. I assume that $f$ is not treated as a constant here. $n$ is the number of elements and $c$ is just some constant. My take is the following, and note that the binomial coefficients is converted into its lower bound:
0003 | $$\binom{n}{2b+1}\left( \frac{2}{cn\cdot 2^f}\right)^{2b}= \frac{n^{2b+1}}{(2b+1)^{2b+1}} \frac{(2b+1)^{2b+1}}{2^f\cdot cn} = \frac{n}{1}\cdot\frac{1}{4^{f}\cdot n}=\frac{1}{4f}$$ Note that $n^{2b+1}=n^{2b}\cdot n = 2n = n$ because its algorithms and their analysis.. My result above is not the same as the one at the top of this question. How does the author derive that? How is $n$ kept in his result?
\end{Vsrc}
\textbf{A YES means:} the proof verification document faithfully represents both the source context and the challenged claim, without dropping a condition or changing its meaning.
\section*{Question 2 --- Does the proposed counterexample disprove that claim?}
\textbf{Claim to test (repeated for convenience).}
\begin{Vsrc}
$$\binom{n}{2b+1}\left( \frac{2}{cn\cdot 2^f}\right)^{2b}= \frac{n^{2b+1}}{(2b+1)^{2b+1}} \frac{(2b+1)^{2b+1}}{2^f\cdot cn} = \frac{n}{1}\cdot\frac{1}{4^{f}\cdot n}=\frac{1}{4f}$$ Note that $n^{2b+1}=n^{2b}\cdot n = 2n = n$
\end{Vsrc}
\subsection*{How the record says the argument uses this claim}
This is the record's description, not an established fact; check it against the source.
\begin{Vsrc}
The submitted derivation chain claims that binom(n,2b+1)*(2/(c*n*2^f))^(2b) equals (after converting the binomial to its lower bound (n/(2b+1))^(2b+1)) the expression n^(2b+1)/(2b+1)^(2b+1) * (2b+1)^(2b+1)/(2^f*c*n), and ultimately equals 1/(4f), using the step n^(2b+1) = n^(2b)*n = 2n = n. Each of these equalities (even read up to constant factors, as the author intends) is false.
\end{Vsrc}
\subsection*{Proposed counterexample (exact values or construction)}
\begin{Vsrc}
Take n = 5, b = 1, c = 1, f = 2 (admissible: n >= 2b+1 = 3, b,f >= 1, c > 0). The four quantities the proof chains together with equals signs are, in exact rational arithmetic:
  binom(5,3)*(2/(c*n*2^f))^(2b) = 10*(2/20)^2 = 1/10                        (true LHS)
  (5/3)^3*(2/20)^2 = 5/108                                                  (LHS after the author's own substitution of the binomial lower bound)
  5^3/3^3 * 3^3/(2^2*1*5) = 25/4                                            (first RHS)
  (5/1)*1/(4^2*5) = 1/16                                                    (second RHS)
  1/(4*2) = 1/8                                                             (final RHS)
Claim says 1/10 = 25/4 = 1/16 = 1/8; in fact these are pairwise distinct, with left/right ratios 2/125, 100 and 1/2. The side remark also fails here: n^(2b+1) = 125, 2n = 10, n = 5, so 125 != 10 != 5.
Smallest witness in enumeration order: n = 3, b = 1, c = 1/8, f = 1, where LHS = 64/9 but the first RHS = 36 (ratio 16/81) and the second RHS = 1/4 (ratio 144).
Exhaustively, across all 96720 enumerated assignments the whole chain holds at ZERO points, exactly or even within a factor 10. Moreover the refutation is quantified over every positive real constant c and every constant factor K: for two assignments sharing (b,f), c enters both sides of a link with the same exponent and cancels from the quotient of the two discrepancy ratios, so that quotient is c-free. At (n,b,f) = (9,4,1) vs (160,4,1) the c-free quotient is 1152921504606846976000000000000000/1744146729569927314733951421 ~ 6.6e5 for link 1, 2684354560000000/4782969 ~ 5.6e8 for the lower-bound variant, and 429496729600000000/43046721 ~ 1.0e10 for link 2 (exactly (n/9)^(2b), hence unbounded in n); for the last link, f = 1 vs f = 12 gives 1048576/3 ~ 3.5e5. Since the link can hold within factor K at both points only if K^2 >= that quotient, no c and no constant K repairs any of them.
\end{Vsrc}
\subsection*{Why the record says the counterexample works}
Check this explanation independently; it is evidence to assess, not an answer key.
\begin{Vsrc}
The proof asserts the chain binom(n,2b+1)*(2/(c*n*2^f))^(2b) = n^(2b+1)/(2b+1)^(2b+1) * (2b+1)^(2b+1)/(2^f*c*n) = (n/1)*1/(4^f*n) = 1/(4f), justified by "n^(2b+1) = n^(2b)*n = 2n = n". This is a finite, parameter-wise refutable statement: each link is an identity in the integers n, b, f and the positive rational c, so a single admissible assignment where two sides differ refutes it.

The script evaluates all five links in exact Fraction arithmetic over n in 3..160, b in 1..4 (with 2b+1 <= n), f in 1..12 and c over a 13-point positive-rational grid (96720 assignments, no floats in any decision, no randomness, no early exit). Both readings of the left-hand side are tested: the genuine binomial coefficient and the author's own replacement of it by the lower bound (n/(2b+1))^(2b+1), so the refutation does not depend on which is meant.

Findings. (i) The whole chain holds at zero assignments, exactly or within a factor 10. (ii) Link 1 (both readings) holds exactly nowhere: the left side carries the factor (2/(c*n*2^f)) to the power 2b, whereas the right side carries (2^f*c*n) to the first power only, and the 2^(2b) numerator vanishes; algebraically the two sides are n^(2b)/(2^f*c) versus binom(n,2b+1)*2^(2b)/(c^(2b)*n^(2b)*2^(2bf)), which are different rational functions. (iii) Link 2 reduces to n^(2b)/(2^f*c) = 1/4^f, i.e. c = n^(2b)*2^f, so it can be satisfied only by letting the "constant" c depend on n and b; in the grid it is hit at exactly one point, (n,b,c,f) = (4,1,32,1). (iv) The c-free ratio-quotient argument (part C of the script) shows the required c grows like n^(2b), so no fixed c and no fixed constant factor K can hold links 1, 1', 2 or 3 simultaneously at two different values of n (or, for link 3, of f). (v) The side remark is false everywhere and for structural reasons: n^(2b)*n = 2n forces n^(2b) = 2, impossible for integers n >= 2, b >= 1 (where n^(2b) >= 4) and false at n = 1; and 2n = n forces n = 0, excluded by n >= 2b+1 >= 3. Exhaustive search finds 0 assignments with n^(2b) = 2 and 0 with 2n = n.

Honest caveats reported by the script rather than suppressed: the FINAL link 1/4^f = 1/(4f) is genuinely TRUE at f = 1 (4 = 4*1) and false for every f >= 2, so the target claim's blanket "each of these equalities is false" is marginally too strong at that one link for that one value of f; and link 2 has the single coincidental grid hit noted above. Neither rescues the chain: at f = 1 links 1, 1' and 2 still fail (see the n = 3, b = 1, c = 1/8, f = 1 witness, where 64/9 != 36 != 1/4), and the coincidental link-2 hit requires a non-constant c.
\end{Vsrc}
\subsection*{How to decide}
\begin{enumerate}
\item List every assumption that applies to the challenged claim.
\item Check that the proposed example satisfies every one of those assumptions.
\item Recompute the claimed outcome and the actual outcome using exact reasoning.
\item Confirm that the conflict is with the challenged claim itself, not an unrelated sentence.
\end{enumerate}
\textbf{A YES means:} the example is fully specified, satisfies all relevant assumptions, and makes the challenged claim false. If a necessary value, condition, or calculation is missing, choose CANNOT DETERMINE.
\section*{Question 3 --- Does the final Lean theorem statement verify this same counterexample?}
\textbf{What to decide.} A Lean file can compile while proving the wrong theorem. Compare these three links independently:
\begin{enumerate}
\item The Lean definitions represent the same mathematical objects, assumptions, and operations as the challenged claim.
\item The Lean construction uses the same concrete values or example shown in Question 2.
\item The final theorem statement asserts that the challenged claim fails on that example, not merely a weaker or nearby statement.
\end{enumerate}
\subsection*{Selected final Lean theorem statement (verbatim UTF-8 excerpt)}
This is an exact contiguous slice of the certificate file, from the theorem keyword through the original \texttt{:=} token where its proof begins. It is shown first so the formal target is not buried in the complete file.
\begin{Vlean}
theorem SE_math_4219732__q_refuted :
    decide (5 ≥ 2 * 1 + 1 ∧ 1 ≥ 1 ∧ 2 ≥ 1) = true
    ∧ fracEq trueLhs ⟨1, 10⟩ = true
    ∧ fracEq substitutedLhs ⟨5, 108⟩ = true
    ∧ fracEq firstRhs ⟨25, 4⟩ = true
    ∧ fracEq secondRhs ⟨1, 16⟩ = true
    ∧ fracEq finalRhs ⟨1, 8⟩ = true
    ∧ printedChain = false :=
\end{Vlean}
\subsection*{Project-local imported Lean files (complete, verbatim UTF-8)}
This section follows the certificate's import commands recursively and prints every repository-owned Lean source file on which the final theorem statement depends. Each file is copied in full, including its definitions and controls. External Mathlib, Std, and Lean Init source is not duplicated.
No project-local imported Lean files were found. External Mathlib, Std, and Lean Init source is not duplicated in this proof verification document.
\subsection*{Complete counterexample Lean file (verbatim UTF-8)}
The complete Lean file is included below so you can check every definition and supporting result used by the final theorem statement. Use it together with the project-local imported Lean files above.
\textbf{Lean certificate file:}
\begin{Vsrc}
partB_feasibility/se_q_certificates/SE_math_4219732__q.lean
\end{Vsrc}
\begin{Vlean}
/-!
# SE_math_4219732__q — the displayed chain of rational equalities fails

* Source: https://math.stackexchange.com/questions/4219732
* Route: `SELF_CONTAINED_LEAN` (core Lean 4, no imports)

## Claim component modeled

The source joins several parameterized rational expressions by equality.  At
the admissible assignment `(n,b,c,f) = (5,1,1,2)`, their exact values are
`1/10`, `5/108`, `25/4`, `1/16`, and `1/8`; hence the chain is false.

## Modelling judgments

Fractions are stored as numerator/positive denominator pairs.  Equality is
cross multiplication, so no division, floating point, normalization axiom, or
library is trusted.  The expressions below reproduce the arithmetic before
reduction.  The negative control checks two different representations of the
same fraction, ensuring that `fracEq` is not merely rejecting every pair.
-/

set_option maxRecDepth 10000
set_option maxHeartbeats 4000000

namespace SE_math_4219732_q

structure Frac where
  num : Nat
  den : Nat
  deriving DecidableEq

def fracEq (a b : Frac) : Bool := decide (a.den > 0 ∧ b.den > 0 ∧ a.num * b.den = b.num * a.den)

def mul (a b : Frac) : Frac := ⟨a.num * b.num, a.den * b.den⟩

def pow : Frac → Nat → Frac
  | _, 0 => ⟨1, 1⟩
  | a, k + 1 => mul a (pow a k)

def trueLhs : Frac := mul ⟨10, 1⟩ (pow ⟨2, 20⟩ 2)
def substitutedLhs : Frac := mul (pow ⟨5, 3⟩ 3) (pow ⟨2, 20⟩ 2)
def firstRhs : Frac := mul ⟨5^3, 3^3⟩ ⟨3^3, 2^2 * 1 * 5⟩
def secondRhs : Frac := mul ⟨5, 1⟩ ⟨1, 4^2 * 5⟩
def finalRhs : Frac := ⟨1, 4 * 2⟩

def printedChain : Bool :=
  fracEq trueLhs firstRhs
  && fracEq firstRhs secondRhs
  && fracEq secondRhs finalRhs

/-- **Main certificate.**  All parameters are admissible, but the equality chain fails. -/
theorem SE_math_4219732__q_refuted :
    decide (5 ≥ 2 * 1 + 1 ∧ 1 ≥ 1 ∧ 2 ≥ 1) = true
    ∧ fracEq trueLhs ⟨1, 10⟩ = true
    ∧ fracEq substitutedLhs ⟨5, 108⟩ = true
    ∧ fracEq firstRhs ⟨25, 4⟩ = true
    ∧ fracEq secondRhs ⟨1, 16⟩ = true
    ∧ fracEq finalRhs ⟨1, 8⟩ = true
    ∧ printedChain = false := by decide

/-- Negative control: cross multiplication accepts equivalent, unreduced fractions. -/
theorem negative_control_fraction_equality :
    fracEq ⟨1, 10⟩ ⟨2, 20⟩ = true := by decide

end SE_math_4219732_q

\end{Vlean}
\subsection*{Proof-safety information}
The automated scan found no unfinished proofs, new axiom declarations, unsafe declarations, external code hooks, or uses of Lean's compiled decision procedure. This scan concerns proof construction only; it does not establish that the theorem matches the source claim.
\textit{Technical scan counts:} native\_decide=0, sorry=0, axiom\_decl=0, unsafe=0, extern=0. \textit{Technical tier:} KERNEL-ONLY.
\textbf{A YES means:} Lean represents the same objects and assumptions, uses the same proposed counterexample, and the final theorem statement asserts the same contradiction. Successful compilation by itself is not enough.
\clearpage
\section*{Human sign-off}
Complete this block from your own review of Questions 1--3. Do not use the automated appendix below as an answer key.\par
\textbf{Reviewer name:} \rule{9cm}{0.4pt}\par
\textbf{Date:} \rule{4cm}{0.4pt}\par
\medskip\textbf{1. Does the proof verification document accurately represent the source and challenged claim?}\par
\answerchoices\par
\decisionexplanation
\medskip\textbf{2. Does the proposed counterexample satisfy the assumptions and make that claim false?}\par
\answerchoices\par
\decisionexplanation
\medskip\textbf{3. Does the final Lean theorem statement represent and verify that same counterexample?}\par
\answerchoices\par
\decisionexplanation
\medskip\textbf{Overall decision:}\par
\checkbox\ VALIDATED (all three answers are YES)\par
\checkbox\ REJECTED (at least one answer is NO)\par
\checkbox\ NEEDS ANOTHER REVIEWER (at least one CANNOT DETERMINE and no NO)\par
\medskip\textbf{Additional comments.} Use this space for any extra source, counterexample, calculation, or Lean details that support your decisions.\par\rule{\linewidth}{0.4pt}\par\vspace{7pt}\rule{\linewidth}{0.4pt}
\clearpage
\section*{Appendix --- Prior automated review (background only)}
\textbf{These results are not an answer key.} They come from the earlier Opus/Fable verification campaign and are printed only so an auditor can see what was previously checked. The human decisions above must be made independently. Raw status codes are retained in small type for traceability.
\textbf{Overall automated disposition:} Cleared after a follow-up reviewer resolved an automated concern.
\par\textit{Audit code:} \nolinkurl{PASS_ADJUDICATED}
\subsection*{Independent search with no access to the stored example}
\textbf{Result:} An independent search found a counterexample without seeing the stored example.
\par\textit{Audit code:} \nolinkurl{WITNESS_FOUND}
\par\textbf{Cases checked:}
\begin{Vsrc}
Full enumeration of the grid n in 0..14, b in 0..3 (2b+1 in {1,3,5,7}), c in 1..5, f in 1..5 = 1500 tuples. Each of the five equality links A-E was tested independently on every tuple where both sides are defined (tuples with a zero denominator, i.e. n=0, were skipped for the affected link, not counted as failures). 1400 tuples defined for links A/B/C; 1500 for links D/E.
\end{Vsrc}
\subsection*{Independent re-check of the proposed counterexample}
\textbf{Result:} A separate checker recomputed the proposed counterexample and confirmed it.
\par\textit{Audit code:} \nolinkurl{CONFIRMED}
\par\textbf{Checker's note:}
\begin{Vsrc}
Independent from-scratch exact recomputation reproduces all five quantities the record lists for n=5,b=1,c=1,f=2 (1/10, 5/108, 25/4, 1/16, 1/8) and all three stated ratios (2/125, 100, 1/2). The quoted chain of equalities fails at every link on the admissible witness, so the witness genuinely refutes the quoted claim.
\end{Vsrc}
\subsection*{Independent review of the source-to-Lean correspondence}
\textbf{Result:} The Lean proof establishes a related failure, but at least one part of the correspondence needs review.
\par\textit{Audit code:} \nolinkurl{PARTIAL}
\par\textbf{Theorem summary from that reviewer:}
\begin{Vsrc}
theorem SE_math_4219732__q_refuted is a single `by decide` conjunction over a hand-rolled Nat/Nat `Frac` type with cross-multiplication equality `fracEq`. It asserts: (a) admissibility `5 >= 2*1+1 AND 1 >= 1 AND 2 >= 1` (i.e. n >= 2b+1, b >= 1, f >= 1 at n=5, b=1, f=2; c=1 is not range-checked, only used as the literal 1); (Shortened here; the complete note remains in the machine-check JSON.)
\end{Vsrc}
\par\textbf{Difference noted by that reviewer:}
\begin{Vsrc}
Link 2 (witness values) and link 3 (the theorem refutes the quoted chain, not an easier statement) both hold. All five Lean expressions were re-evaluated independently in exact rational arithmetic and reproduce the record's values 1/10, 5/108, 25/4, 1/16, 1/8 exactly, and the record's witness (n,b,c,f) = (5,1,1,2) is recoverable from the encodings: 2*1+1 = 3 and the exponents 2b=2, 2b+1=3 in `pow`, 2^f as 2^2 and 4^f as 4^2, c=1 explicit in firstRhs's denominator 2^2*1*5, and 4f as 4*2. (Shortened here; the complete note remains in the machine-check JSON.)
\end{Vsrc}
\subsection*{Follow-up review of an automated concern}
\textbf{Result:} Lean uses a documented simplification judged not to change the refutation.
\par\textit{Audit code:} \nolinkurl{DOCUMENTED_MODELING_CHOICE}.\par\textit{Confidence:} \nolinkurl{high}.\par\textit{Recommended action:} \nolinkurl{NONE}.
\par\textbf{Conclusion of the follow-up review:}
\begin{Vsrc}
At the admissible assignment (n,b,c,f) = (5,1,1,2) (n >= 2b+1 = 3, b >= 1, f >= 1, c > 0), the five expressions the source joins with equals signs are, in exact rationals: binom(5,3)*(2/(1*5*2^2))^2 = 10*(1/10)^2 = 1/10; the author's own binomial-lower-bound reading (5/3)^3*(1/10)^2 = 5/108; 5^3/3^3 * 3^3/(2^2*1*5) = 25/4; (5/1)*1/(4^2*5) = 1/16; and 1/(4*2) = 1/8. These are pairwise distinct, with link ratios 2/125 (and 1/135 for the lower-bound reading), 100 and 1/2, so no link holds even within a factor of 10. (Shortened here; the complete note remains in the machine-check JSON.)
\end{Vsrc}
\textit{This appendix is only background. You do not need any separate machine-output files to answer the three questions.}
\end{document}